% ==============================================================================
% Plantilla Institucional de Trabajo Terminal — UPIIZ IPN
% Archivo: capitulos/protocolo/05-estado-del-arte.tex
% ==============================================================================
% ESTRUCTURA NORMATIVA DE PROTOCOLO: CAPÍTULO 5 - ESTADO DEL ARTE
% ==============================================================================
% LÍMITE INSTITUCIONAL: Máximo 4 páginas.
%
% Requisitos normativos (Anexo 1):
% - Presentar la situación actual de la investigación técnica, científica e industrial.
% - Describir desarrollos innovadores o recientes directamente relacionados con el tema.
% - Redactar mediante paráfrasis formal y referencias bibliográficas correctas en formato IEEE.
%
% Buena práctica sugerida (no obligatoria):
% - Sintetizar las soluciones o enfoques en una tabla comparativa opcional.
% ==============================================================================

\chapter{Estado del Arte}
\label{cap:proto-estado-del-arte}

% [INSTRUCCIÓN PARA EL ALUMNO]:
% Desarrolle en un máximo de 4 páginas la revisión del estado actual de la tecnología,
% investigación y desarrollo industrial relacionados con el problema de su proyecto.

\section{Situación actual de la investigación técnica, científica e industrial}

[Describa el panorama contemporáneo de la técnica y la ciencia en el campo del proyecto, citando fuentes formales (artículos científicos, patentes, desarrollos tecnológicos comerciales o normas recientes)].

\section{Desarrollos innovadores y tecnologías recientes}

[Analice las soluciones innovadoras y productos afines reportados recientemente en la literatura especializada que guarden relación directa con el problema abordado].

% ------------------------------------------------------------------------------
% SUGERENCIA EDITORIAL (BUENA PRÁCTICA RECOMENDADA, NO REQUISITO OBLIGATORIO):
% Si el alumno desea sintetizar las tecnologías revisadas, puede incluir una tabla comparativa:
%
% \begin{table}[htbp]
%     \centering
%     \caption{Comparación técnica de enfoques recientes en el estado del arte.}
%     \label{tab:proto-comparativa-opcional}
%     \begin{tabularx}{\textwidth}{p{3.5cm} X X}
%         \toprule
%         \textbf{Desarrollo / Enfoque} & \textbf{Características destacadas} & \textbf{Limitaciones observadas} \\
%         \midrule
%         [Solución A] & [Características] & [Limitaciones] \\
%         [Solución B] & [Características] & [Limitaciones] \\
%         \bottomrule
%     \end{tabularx}
% \end{table}
